\documentclass[10pt,a4paper]{amsart}
\usepackage[margin=19mm]{geometry}
\usepackage{amsmath,amssymb,mathtools}
\usepackage[scr=rsfs,cal=euler]{mathalfa}
\usepackage{xfrac}
\usepackage[unicode,hidelinks,bookmarks=false]{hyperref}
\newcommand{\ie}{i.e.\ }
\newcommand{\cf}{cf.\ }
\newcommand{\resp}{respectively\ }
\newcommand{\Subsection}{\subsection}
\providecommand{\nobreakdash}{\nobreak\hbox{-}}
\setlength{\emergencystretch}{2em}
\multlinegap0pt
\usepackage{amssymb}
\usepackage{tikz}
\newtheorem{theorem}{Theorem}
\newcommand{\N}{{\mathbb N}}
\newcommand{\longto}{\longrightarrow}
\newcommand{\quilt}{\operatorname{Quilts}}
\newcommand{\rank}{\operatorname{rank}}
\DeclareMathOperator{\rk}{quiltrank}
\title{Generalized Rank Functions \\ and Quilts of Alternating Sign Matrices}
\author{Sara Billey, Matja\v{z} Konvalinka}
\date{Source: arXiv:2412.03236v2 (CC BY 4.0)}
\begin{document}
\maketitle

\noindent\emph{Source and licence.} Generalized Rank Functions \\ and Quilts of Alternating Sign Matrices. Version arXiv:2412.03236v2 (CC BY 4.0), distributed by arXiv under the Creative Commons Attribution 4.0 International licence, http://creativecommons.org/licenses/by/4.0/, as recorded in the arXiv licence metadata for this exact version and shown on the abstract page https://arxiv.org/abs/2412.03236v2. Full licence notice: \texttt{licenses/arxiv-2412.03236-CC-BY-4.0.txt}.\par\smallskip

\noindent\textit{Editorial scope.} Counted unit: the article's theorem thm:lattice (source0.tex lines 838--846). The article's complete original proof of it is delivered with it (source0.tex lines 865--928, one \texttt{proof} environment of 64 lines, which the article places away from the statement). No mathematics is written, repaired or re-derived: the statement and the proof are the article's own lines, in the article's own notation, and no retained line is substituted or re-flowed.  No further source-local context is needed: every cross-reference of every retained line resolves inside the two delivered spans.  Nothing was omitted from any retained span, so the package carries no omission record.  No retained line contains a citation, so the wrapper carries no bibliography and no bibliography entry.  No retained line contains a figure environment, an image-inclusion command or drawing code, so no image is delivered and the package declares no asset dependency.  Editorial formatting only, in addition: an AMS \texttt{amsart} preamble that restates the article's own declarations and macros used by the retained lines, namely the environments \texttt{theorem} on separate counters, together with the standard packages those lines need.  The retained environments print in this document as Theorem 1 in order of appearance, whereas the article prints the same items with the numbering of its own full document.  The complete native source of the article as distributed in the arXiv e-print for this exact version is bundled unchanged as \texttt{sources/169603/source0.tex}.
\begin{theorem} \label{thm:lattice}
Let $P,Q$ be finite ranked posets with least and greatest elements.
The poset $\quilt(P,Q)$ is a distributive lattice ranked by $$\rk f =
\sum_{x \in P,\, y \in Q} f(x,y) - \sum_{x \in P,\, y \in Q} f_{\hat
0}(x,y),$$ where $f_{\hat 0}(x,y) = \max\{0,\rank x + \rank y -
\max\{\rank P,\rank Q\}\}$ is the least element of $\quilt(P,Q)$. The greatest
element of $\quilt(P,Q)$ is $f_{\hat 1}(x,y) = \min\{\rank x,\ \rank
y\}$.
\end{theorem}

\begin{proof}[Proof of Theorem~\ref{thm:lattice}] To prove that $\quilt(P,Q)$ is a lattice, observe that for every $f,g \in
\quilt(P,Q)$, the pointwise minimum and maximum of $f$ and $g$ is
again in $\quilt(P,Q)$ by definition.  These quilts are $f\wedge g$ and $f \vee g$
respectively, hence $\quilt(P,Q)$ is a lattice. Distributivity follows
from the fact that $$\max\{a,\min\{b,c\}\} =
\min\{\max\{a,b\},\max\{a,c\}\}$$ and $$\min\{a,\max\{b,c\}\} =
\max\{\min\{a,b\},\min\{a,c\}\}$$ hold for $a,b,c \in \N$.


It is easy to check that $f_{\hat 0}$ and $f_{\hat 1}$ are in
$\quilt(P,Q)$ by checking the properties of the definition.  To see
$f_{\hat 1}$ is the unique maximal element in $\quilt(P,Q)$, consider
any $(x,y) \in P \times Q$ with $\rank x = r$, $\rank y = s$. Since
$P,Q$ are ranked posets, there exist maximal chains
\[
\hat 0_P = x_0 \lessdot x_1 \lessdot \dots
\lessdot x_r = x \lessdot \dots \lessdot x_k = \hat 1_P
\]
and
\[
\hat 0_Q = y_0 \lessdot y_1 \lessdot \dots \lessdot y_s = y \lessdot
\dots \lessdot y_n = \hat 1_Q.
\]
Then, for any $f \in \quilt(P,Q)$, we know
$f(x,y) \leq f(x_{r-1},y) + 1 \leq \ldots \leq f(\hat 0_P,y) + r = r$
and $f(x,y) \leq f(x,y_{s-1}) + 1 \leq \ldots \leq f(x,\hat 0_Q) + s =
s$, so $f(x,y) \leq \min\{r,s\} = f_{\hat 1}(x,y)$. Hence, $f_{\hat 1}$ is maximal.
On the other hand, $f(x,y) \geq f(x,y_{s+1}) - 1 \geq \dots \geq
f(x,\hat 1_Q) - (\rank Q-s) \geq f(x_{r+1},\hat 1_Q) - 1 - (\rank Q-s) \geq \dots
\geq f(\hat 1_P,\hat 1_Q) - (\rank P-r) - (\rank Q-s) = \min\{\rank P,\rank Q\} - \rank P - \rank Q + r + s = r + s - \max\{\rank P,\rank Q\}$ and so $f(x,y) \geq f_{\hat 0}(x,y)$.
Therefore, $f_{\hat 0}$ is minimal in $\quilt(P,Q)$.



To see $\quilt(P,Q)$ is ranked by the given function, choose $f,g \in
\quilt(P,Q)$ with $f < g$. Our goal is to prove that $f \lessdot g$ if and
only if $f$ and $g$ differ in exactly one $(x,y) \in P \times  Q$, and $g(x,y) =
f(x,y) + 1$. The condition is clearly sufficient, let us prove that it is also necessary.


Since $f<g$, the directed graph $G_\Delta(f,g)$ constructed above is
non-empty. We claim that it has no directed cycles, and that it
therefore has at least one sink.  To observe the claim, note that if
$(x,y) \to (x',y')$, then $f(x,y) \geq f(x',y')$, and if $f(x,y) =
f(x',y')$, then $\rank(x, y) < \rank(x', y')$ in $P \times Q$. If $(x,y)
\to (x',y') \to (x'',y'') \to \dots \to (x,y)$, then the value of $f$
must stay constant (if it strictly decreases, it can never increase to
$f(x,y)$ again), and that implies that $\rank(x,y) < \rank(x',y') <
\rank(x'',y'') < \dots < \rank(x,y)$, which is a contradiction.

Take $(x,y)$ to be an arbitrary sink in $G_\Delta(f,g)$.  Observe by
choice of $(x,y)$ that if $(x,y) \lessdot (x',y')$, then $f(x',y') =
f(x,y) + 1$, and if $(x',y') \lessdot (x,y)$, then $f(x',y') =
f(x,y)$. Since $f(\hat 0_{P},\hat 0_{Q})=g(\hat 0_{P},\hat 0_{Q})=0$
and $f(\hat 1_{P},\hat 1_{Q})=g(\hat 1_{P},\hat 1_{Q})=\min\{\rank P,\
\rank Q\}$ by definition of a quilt, $(x,y)$ is not $(\hat 0_{P},\hat
0_{Q})$ or $(\hat 1_{P},\hat 1_{Q})$.  Therefore, the function $f'
\colon P \times Q \longto \N$, defined to be equal to $f$ everywhere
except $f'(x,y)=f(x,y)+1$, is also a quilt of type $(P,Q)$, hence
$f\lessdot f'$.  Furthermore, since $(x,y)$ is a vertex in
$G_\Delta(f,g)$, we know $f'(x,y)=f(x,y)+1\leq g(x,y)$, so $f'\leq g$
as well. Therefore, we conclude that $f \lessdot g$ if and only if
$f'=g$.
\end{proof}

\end{document}
